\section{Experiments}
\label{sec:experiments}

\newcommand{\figplaceholder}[1]{%
  \fbox{\parbox[c][0.24\textheight][c]{0.92\linewidth}{\centering\small #1}}}

The protocol compares estimated curves against simulated curves on production
traces, at capacities the model never saw during fitting. Every number in this
section is a placeholder pending the runs described in
\cref{sec:appendix-repro}.

\subsection{Protocol}

Each trace is split by time into a fit segment, a validation segment, and a
held-out segment, in that order, with no interleaving. Reuse distances are
computed once over the whole trace so that distances crossing a split boundary
are not truncated. Models are fitted on the fit segment, hyperparameters and
the censoring correction are chosen on the validation segment, and all reported
errors are on the held-out segment. Ground truth is simulated on the held-out
segment only. This ordering prevents the model from observing any reference
whose distance it is asked to predict.

\subsection{Capacity Grid}

Capacities are fixed before any fitting, so that the same grid is used for
simulation, for evaluation, and for every model.

\begin{table}[ht]
\centering
\caption{Capacity grid used for simulation and evaluation.}
\label{tab:capacity-grid}
\begin{tabular}{@{}S[table-format=2.0] l S[table-format=4.0]@{}}
\toprule
{$C$ (GiB)} & Regime & {Blocks ($\times 10^3$)} \\
\midrule
1  & below knee      & \ldots \\
2  & below knee      & \ldots \\
4  & near knee       & \ldots \\
8  & near knee       & \ldots \\
16 & above knee      & \ldots \\
32 & above knee      & \ldots \\
64 & saturation      & \ldots \\
\bottomrule
\end{tabular}
\end{table}

\subsection{Clusters}

Traces come from the 2020 production cache analysis of \cite{yang2020}. Only
clusters whose traces are complete over the fit and held-out windows are used,
and the selection rule is applied before any model is fitted.

\begin{table}[ht]
\centering
\caption{Clusters under study. Values are placeholders.}
\label{tab:clusters}
\begin{tabular}{@{}l S[table-format=5.1] S[table-format=4.1] S[table-format=2.2] l@{}}
\toprule
Cluster & {Requests ($\times 10^6$)} & {Distinct keys ($\times 10^3$)} & {Mean object (KiB)} & Workload class \\
\midrule
C1 & \ldots & \ldots & \ldots & placeholder \\
C2 & \ldots & \ldots & \ldots & placeholder \\
C3 & \ldots & \ldots & \ldots & placeholder \\
C4 & \ldots & \ldots & \ldots & placeholder \\
\bottomrule
\end{tabular}
\end{table}

\subsection{Results}

\Cref{tab:mrc-error} reports curve error against simulation, per cluster, for
both hazard parametrizations. \Cref{tab:calibration} reports hazard calibration
and the integrated Brier score. Neither table carries a value until the runs
complete.

\begin{table}[ht]
\centering
\caption{MRC error against simulation on the held-out segment. Lower is better.}
\label{tab:mrc-error}
\begin{tabular}{@{}l S[table-format=1.3] S[table-format=1.3] S[table-format=1.3] S[table-format=1.3]@{}}
\toprule
& \multicolumn{2}{c}{R-as-feature} & \multicolumn{2}{c}{Per-bin} \\
\cmidrule(lr){2-3}\cmidrule(lr){4-5}
Cluster & {Max} & {Mean} & {Max} & {Mean} \\
\midrule
C1 & \ldots & \ldots & \ldots & \ldots \\
C2 & \ldots & \ldots & \ldots & \ldots \\
C3 & \ldots & \ldots & \ldots & \ldots \\
C4 & \ldots & \ldots & \ldots & \ldots \\
\bottomrule
\end{tabular}
\end{table}

\begin{table}[ht]
\centering
\caption{Hazard calibration and integrated Brier score on the held-out segment.}
\label{tab:calibration}
\begin{tabular}{@{}l S[table-format=1.3] S[table-format=1.4] S[table-format=1.4]@{}}
\toprule
Cluster & {ECE} & {Brier} & {IBS} \\
\midrule
C1 & \ldots & \ldots & \ldots \\
C2 & \ldots & \ldots & \ldots \\
C3 & \ldots & \ldots & \ldots \\
C4 & \ldots & \ldots & \ldots \\
\bottomrule
\end{tabular}
\end{table}

\begin{figure}[ht]
\centering
\IfFileExists{../artifacts/figures/mrc-curves.pdf}{%
  \includegraphics[width=\linewidth]{../artifacts/figures/mrc-curves.pdf}%
}{\figplaceholder{PLACEHOLDER: estimated and simulated MRCs per cluster, generated by \texttt{src/mrc/plots/}}}
\caption{Estimated and simulated miss-ratio curves. Simulation is the solid
line, the two parametrizations are dashed.}
\label{fig:mrc-curves}
\end{figure}

\begin{figure}[ht]
\centering
\IfFileExists{../artifacts/figures/calibration.pdf}{%
  \includegraphics[width=0.7\linewidth]{../artifacts/figures/calibration.pdf}%
}{\figplaceholder{PLACEHOLDER: reliability diagram for the hazard model, generated by \texttt{src/mrc/plots/}}}
\caption{Reliability diagram for the discrete hazard, pooled over bins.}
\label{fig:calibration}
\end{figure}

\begin{figure}[ht]
\centering
\IfFileExists{../artifacts/figures/cluster-error.pdf}{%
  \includegraphics[width=\linewidth]{../artifacts/figures/cluster-error.pdf}%
}{\figplaceholder{PLACEHOLDER: per-cluster error against capacity, generated by \texttt{src/mrc/plots/}}}
\caption{Curve error against capacity, per cluster, for the per-bin
parametrization.}
\label{fig:cluster-error}
\end{figure}
